Hysterezní smyčka $P(E)$ ukazuje, že polarizace závisí i na předchozím elektrickém poli:
\begin{itemize}
\item Rostoucí pole přepólovává domény až k nasycení.
\item Při návratu na $E=0$ zůstává \textbf{remanentní polarizace $P_r$}.
\item Opačné \textbf{koercitivní pole $-E_c$} sníží polarizaci na nulu; další zesílení pole vede k zápornému nasycení.
\item Zpětná větev prochází zápornou remanencí $-P_r$ a kladným koercitivním polem $+E_c$, čímž se smyčka uzavře.
\end{itemize}
\textbf{Curieova teplota $T_C$} je teplota přechodu z feroelektrického do paraelektrického stavu: zaniká spontánní polarizace. U BaTiO$_3$ je přibližně $120\,^{\circ}\mathrm C$ a v jejím okolí dosahuje permitivita maxima.
